\documentclass[10pt,a4paper]{amsart}
\usepackage[margin=19mm]{geometry}
\usepackage{amsmath,amssymb,mathtools}
\usepackage[scr=rsfs,cal=euler]{mathalfa}
\usepackage{xfrac}
\usepackage[unicode,hidelinks,bookmarks=false]{hyperref}
\newcommand{\ie}{i.e.\ }
\newcommand{\cf}{cf.\ }
\newcommand{\resp}{respectively\ }
\newcommand{\Subsection}{\subsection}
\providecommand{\nobreakdash}{\nobreak\hbox{-}}
\setlength{\emergencystretch}{2em}
\multlinegap0pt
\usepackage{amssymb}
\usepackage{enumerate}
\usepackage{mathtools}
\newtheorem{lemma}{Lemma}
	\newcommand{\la}{\lambda}
	\newcommand{\ls}{\lesssim}
\title{Maximal growth of the Stein-Wainger oscillatory integral}
\author{Cheng Zhang, Zhifei Zhu}
\date{Source: arXiv:2603.23238v1 (CC BY 4.0)}
\begin{document}
\maketitle

\noindent\emph{Source and licence.} Maximal growth of the Stein-Wainger oscillatory integral. Version arXiv:2603.23238v1 (CC BY 4.0), distributed by arXiv under the Creative Commons Attribution 4.0 International licence, http://creativecommons.org/licenses/by/4.0/, as recorded in the arXiv licence metadata for this exact version and shown on the abstract page https://arxiv.org/abs/2603.23238v1. Full licence notice: \texttt{licenses/arxiv-2603.23238-CC-BY-4.0.txt}.\par\smallskip

\noindent\textit{Editorial scope.} Counted unit: the article's lemma finitetype (source0.tex lines 280--282). The article's complete original proof of it is delivered with it (source0.tex lines 283--373, one \texttt{proof} environment of 91 lines). No mathematics is written, repaired or re-derived: the statement and the proof are the article's own lines, in the article's own notation, and no retained line is substituted or re-flowed.  No further source-local context is needed: every cross-reference of every retained line resolves inside the two delivered spans.  Nothing was omitted from any retained span, so the package carries no omission record.  No retained line contains a citation, so the wrapper carries no bibliography and no bibliography entry.  No retained line contains a figure environment, an image-inclusion command or drawing code, so no image is delivered and the package declares no asset dependency.  Editorial formatting only, in addition: an AMS \texttt{amsart} preamble that restates the article's own declarations and macros used by the retained lines, namely the environments \texttt{lemma} on separate counters, together with the standard packages those lines need.  The retained environments print in this document as Lemma 1 in order of appearance, whereas the article prints the same items with the numbering of its own full document.  The complete native source of the article as distributed in the arXiv e-print for this exact version is bundled unchanged as \texttt{sources/197064/source0.tex}.
\begin{lemma}\label{finitetype}
	Suppose that $\phi$ has a finite order of vanishing at 0. Then $m(\la)=O(1)$.
\end{lemma}

\begin{proof}
	Assume $\phi\not\equiv 0$ and let $k$ be the smallest odd integer with $\phi^{(k)}(0)\neq 0$.
	Then Taylor's theorem gives
	\[
	\phi(t)=c\,t^k+O(t^{k+2}),\qquad c=\frac{\phi^{(k)}(0)}{k!}\neq 0.
	\]
	Hence there exists $\delta\in(0,1)$ and constants $c_0,C_0>0$ such that for $0<t\le\delta$,
	\begin{equation}\label{eq:finite-type}
		c_0 t^k\le |\phi(t)|\le C_0 t^k,
		\qquad
		|\phi^{(k)}(t)|\ge c_0.
	\end{equation}
	
	Split
	\[
	m(\lambda)=\int_0^\delta \frac{e^{i\lambda\psi(t)}-e^{i\lambda\psi(-t)}}{t}\,dt
	+\int_\delta^1 \frac{e^{i\lambda\psi(t)}-e^{i\lambda\psi(-t)}}{t}\,dt.
	\]
	The second integral is $O(1)$, so it remains to bound the first one, denoted by $m_0(\la)$, uniformly in $\lambda$.
	
	Let $I_j=[2^{-j-1}\delta,2^{-j}\delta]$, $t_j=2^{-j}\delta$ ($j\ge 0$). Dyadically decompose the integral and
	rescale $t=t_js$, $s\in[1/2,1]$. We obtain
	\[
	m_0(\lambda)=\sum_{j\ge 0} J_j(\lambda),
	\qquad
	J_j(\lambda):=\int_{1/2}^1 \frac{e^{i\lambda\psi(t_js)}-e^{i\lambda\psi(-t_js)}}{s}\,ds.
	\]
	Define the scale parameter $\Lambda_j:=\lambda\,t_j^k=\lambda\,\delta^k\,2^{-jk}.$
We claim that \begin{equation}\label{clm}
	|J_j|\ls \min \{\Lambda_j,\ \Lambda_j^{-1/k}\}.
\end{equation} We postpone the proof of this claim and use it to obtain a uniform bound for $m_0(\lambda)$.

 Let $j_*$ be such that $\Lambda_{j_*}\ge 1>\Lambda_{j_*+1}$, i.e.
$2^{j_*}\approx (\lambda\delta^k)^{1/k}$. Then 
\[
\sum_{j\le j_*}|J_j(\lambda)|
\lesssim \sum_{j\le j_*}\Lambda_j^{-1/k}
=\sum_{j\le j_*}(\lambda\delta^k)^{-1/k}2^{j}
\lesssim (\lambda\delta^k)^{-1/k}2^{j_*}
\lesssim 1,
\]
and
\[
\sum_{j> j_*}|J_j(\lambda)|
\lesssim \sum_{j>j_*}\Lambda_j
=\sum_{j>j_*}\lambda\delta^k\,2^{-jk}
\lesssim \lambda\delta^k\,2^{-k(j_*+1)}
\lesssim 1.
\]

Now it remains to prove the claim \eqref{clm}. On the one hand, using $|e^{ix}-e^{iy}|\le |x-y|$ and \eqref{eq:finite-type}, we obtain
	\[
	\big|e^{i\lambda\psi(t)}-e^{i\lambda\psi(-t)}\big|
	\le \la\,|\psi(t)-\psi(-t)|
	=\la\,|\phi(t)|
	\lesssim \la\,t^k \quad (0<t\le\delta).
	\]
 Hence for $s\in[1/2,1]$,
	\[
	\big|e^{i\lambda\psi(t_js)}-e^{i\lambda\psi(-t_js)}\big|
	\lesssim \la\,(t_js)^k\lesssim \la\,t_j^k=\Lambda_j.
	\]
	Therefore
	\begin{equation}\label{eq:small}
		|J_j(\lambda)|
		\ls \int_{1/2}^1 \frac{\Lambda_j}{s}\,ds
		\lesssim \Lambda_j.
	\end{equation}
	
	
	On the other hand, write $J_j=A_j-B_j$ where
	\[
	A_j(\lambda):=\int_{1/2}^1 \frac{e^{i\lambda\psi(t_js)}}{s}\,ds,
	\qquad
	B_j(\lambda):=\int_{1/2}^1 \frac{e^{i\lambda\psi(-t_js)}}{s}\,ds.
	\]
Let $\Phi_j(s)=\psi(t_js)$. 
	Since $k$ is odd and
\[
\phi^{(k)}(t)=\psi^{(k)}(t)+\psi^{(k)}(-t)=2\psi^{(k)}(0)+O(t),
\]
by \eqref{eq:finite-type} we have $|\psi^{(k)}(u)|\gtrsim 1$ for $|u|\le\delta$
	(up to adjusting $\delta$), hence $|\Phi_j^{(k)}(s)|\gtrsim t_j^k$ on $[1/2,1]$. Then a standard van der Corput estimate for phases with a nonvanishing $k$th derivative on a fixed interval yields
	\[
	|A_j(\lambda)|\lesssim (\lambda t_j^k)^{-1/k}=\Lambda_j^{-1/k}.
	\]
	Similarly $|B_j(\lambda)|\lesssim \Lambda_j^{-1/k}$, hence
	\begin{equation}\label{eq:large}
		|J_j(\lambda)|\le |A_j(\lambda)|+|B_j(\lambda)|\lesssim \Lambda_j^{-1/k}.
	\end{equation}
	This proves the claim \eqref{clm} and completes the proof.\end{proof}

\end{document}
